% Part: methods
% Chapter: proofs
% Section: starting-proofs

\documentclass[../../../include/open-logic-section]{subfiles}

\begin{document}

\olfileid{mth}{prf}{str}

\olsection{నిరూపణను ప్రారంభించడం}

అసలు ఎక్కడ మొదలుపెట్టాలి?

మీకు నిరూపించాల్సిన విషయం ఇవ్వబడింది; కనుక నిరూపణలో చివరికి చేరవలసింది అదే (దాన్ని నిరూపించబోతున్నామని మొదట \emph{ప్రకటించవచ్చు}, కానీ దాన్నే పరికల్పనగా వాడకూడదు). కొత్త కాగితం అడుగున నిరూపించాల్సిన విషయాన్ని రాయండి---ఇలా చేస్తే మీ లక్ష్యం దృష్టి నుంచి తప్పదు.

తరువాత మీరు ఉపయోగించగల పరికల్పనలు కొన్ని ఉండవచ్చు (తరువాతి విభాగంలో నిరూపణ \emph{రకం} గురించి చెప్పినప్పుడు ఇది మరింత స్పష్టమవుతుంది). వాటిని కాగితం పైభాగంలో రాసి, అవి పరికల్పనలని స్పష్టంగా చూపండి (అంటే $p$ను పరికల్పనగా తీసుకుంటే, ``$p$ అని అనుకుందాం'' లేదా ``$p$ అని భావిద్దాం'' అని రాయండి). చివరిగా, ప్రశ్నలోని కొన్ని నిర్వచనాలు తెలుసుకోవాల్సి రావచ్చు. ఒక నిర్దిష్ట నిర్వచనాన్ని వాడమని చెప్పవచ్చు; లేదా పరికల్పనలలోనూ మీరు చేరాల్సిన నిష్కర్షలోనూ వివిధ నిర్వచనాలు ఉండవచ్చు. \emph{వాటిని రాసుకొని, వాటి అర్థం మీకు స్పష్టమైందని నిర్ధారించుకోండి.}

మీ నిరూపణను ఎలా అమర్చాలో ప్రశ్న రూపంపైనా ఆధారపడుతుంది. వాక్య రకాన్ని బట్టి నిరూపణను ఎలా ప్రారంభించాలో తరువాతి విభాగం వివరిస్తుంది.

\end{document}
